\chapter{Toxicology}

\medterm{Activated Charcoal}-- Fine black powder produced by controlled pyrolysis of organic materials, processed to yield a massive micro-porous surface area (500-1500 m2/g). Mechanism: non-specific adsorption of toxins onto carbon surface via Van der Waals forces and hydrophobic interactions. Dose: 1 g toxin per 10 g charcoal (10:1 ratio), typical adult dose 50-100 g. Indications: recent ingestion (1-2 hours) of adsorbable toxins. Contraindications: corrosives, hydrocarbons, lithium, iron, ethanol, methanol. Complications: aspiration, constipation, bowel obstruction.

\medterm{Acetaminophen (Paracetamol) Toxicity}-- Most common cause of acute liver failure in Western countries. Mechanism: toxic metabolite NAPQI depletes glutathione, causing centrilobular hepatic necrosis. Treatment: N-acetylcysteine (NAC) within 8-10 hours of ingestion. Nomogram (Rumack-Matthew) guides treatment based on serum acetaminophen level at 4+ hours. Dosing: loading dose 140 mg/kg, maintenance 70 mg/kg every 4 hours for 17 doses (oral) or continuous IV infusion.

\medterm{Acute Liver Failure}-- Rapid loss of liver function with coagulopathy and encephalopathy in non-cirrhotic liver. Causes: acetaminophen toxicity (most common in West), viral hepatitis, toxins (amanitin, halothane), ischemic hepatitis. Management: N-acetylcysteine (early), supportive care, transplantation assessment.

\medterm{Acute Kidney Injury}-- Sudden decline in renal function. Toxic causes: contrast media, aminoglycosides, cisplatin, ethylene glycol, mycotoxins. Mechanisms: direct tubular toxicity, vascular injury, obstruction (crystals). Management: remove toxin, supportive care, dialysis if severe.

\medterm{Air Contaminants}-- See occupational medicine.

\medterm{Alcohol}-- See ethanol.

\medterm{Alcohol Intoxication}-- Acute effects of ethanol: CNS depression, disinhibition, ataxia, slurred speech, vomiting, respiratory depression (severe). BAC >0.08% impairs driving. Treatment: supportive, airway protection, thiamine (prevent Wernicke), glucose if hypoglycemic.

\medterm{Alcohol Withdrawal}-- See withdrawal syndromes.

\medterm{Aluminum Phosphide}-- Fumigant poison released as phosphine gas upon contact with moisture. Mechanism: inhibits cytochrome c oxidase, cellular hypoxia. Treatment: supportive, cardiac monitoring, magnesium, calcium. No specific antidote.

\medterm{Amelogenesis}-- See dental.

\medterm{Amphetamine}-- See stimulants.

\medterm{Anion Gap}-- Calculated as [Na+] - ([Cl-] + [HCO3-]). Normal: 8-12 mEq/L. Elevated in: MUDPILES (Methanol, Uremia, DKA, Propylene glycol, Iron/Isoniazid, Lactic acidosis, Ethylene glycol, Salicylates). Used to identify toxic alcohols, metabolic acidosis.

\medterm{Antidote}-- Substance that counteracts poison toxicity. Specific: naloxone (opioids), flumazenil (benzodiazepines), N-acetylcysteine (acetaminophen), digoxin immune Fab (digoxin), chelators (EDTA, deferoxamine, dimercaprol), atropine + pralidoxime (organophosphates), methylene blue (methemoglobinemia), cyanide antidote kit (nitrite + thiosulfate). Non-specific: activated charcoal, supportive care.

\medterm{Arterial Blood Gas (ABG)}-- Measures pH, PaCO2, PaO2, HCO3-, base excess. Used to assess acid-base status, oxygenation, ventilation. Normal: pH 7.35-7.45, PaCO2 35-45 mmHg, PaO2 80-100 mmHg, HCO3- 22-26 mEq/L. Interpretation: stepwise approach identifies primary disorder and compensation.

\medterm{Aspirin (Salicylate) Toxicity}-- Mixed respiratory alkalosis and metabolic acidosis. Early: tinnitus, hyperventilation, respiratory alkalosis. Late: metabolic acidosis, fever, seizures, coma. Treatment: activated charcoal, urine alkalinization (sodium bicarbonate), hemodialysis if severe (renal failure, pulmonary edema, severe acidosis, very high levels).

\medterm{Barbiturate}-- See sedative-hypnotics.

\medterm{Benzo(a)pyrene}-- Polycyclic aromatic hydrocarbon from combustion. Carcinogenic (lung, skin cancer). Found in tobacco smoke, charred meat, air pollution.

\medterm{Benzene}-- Solvent, industrial chemical. Chronic exposure: aplastic anemia, leukemia (especially AML). Acute: CNS depression, arrhythmias. Prevention: substitution, ventilation, monitoring.

\medterm{Beta-Blocker Overdose}-- Bradycardia, hypotension, hypoglycemia, bronchospasm (non-selective). Treatment: atropine, glucagon (1-10 mg IV), high-dose insulin euglycemia therapy, calcium, vasopressors, pacing if refractory.

\medterm{Bisphenol A (BPA)}-- Endocrine disruptor in plastics. Estrogenic effects, associated with reproductive disorders, obesity, diabetes. Avoid BPA-containing products, especially for children.

\medterm{Blood Alcohol Concentration (BAC)}-- Measure of ethanol in blood. Legal limit for driving: 0.08% (varies by jurisdiction). Metabolism: zero-order (15-20 mg/dL/hour), liver breakdown. Units: g/dL or mg/dL.

\medterm{Bone Marrow Depression}-- See myelosuppression.

\medterm{Carbon Monoxide (CO) Poisoning}-- Inhibits hemoglobin oxygen binding (affinity 200x oxygen). Causes headache, dizziness, nausea, cherry-red skin, syncope, death. Diagnosis: carboxyhemoglobin level (not routine). Treatment: 100% oxygen, hyperbaric oxygen if severe (neurological symptoms, COHb >25%, pregnancy).

\medterm{Calcium Channel Blocker Overdose}-- Hypotension, bradycardia, hyperglycemia. Treatment: calcium gluconate/chloride, glucagon, high-dose insulin, vasopressors, lipid emulsion therapy, pacing.

\medterm{Carboxyhemoglobin}-- See carbon monoxide poisoning.

\medterm{Carcinogen}-- Substance causing cancer. Types: chemical (benzene, asbestos, aflatoxin), physical (radiation), biological (HPV, HBV, HCV). Mechanisms: DNA damage, mutations, epigenetic changes. IARC classifications: Group 1 (carcinogenic), Group 2A (probably carcinogenic), Group 2B (possibly carcinogenic), Group 3 (not classifiable).

\medterm{Carrier State}-- Asymptomatic infection with pathogen capable of transmission. Examples: typhoid carrier, hepatitis B carrier, meningococcal carrier. Importance: public health control, screening, treatment.

\medterm{Cerebellar Ataxia}-- See neurology.

\medterm{Charcoal}-- See activated charcoal.

\medterm{Chelation Therapy}-- Use of chelating agents to bind metals for excretion. Agents: EDTA (lead), deferoxamine (iron), dimercaprol (arsenic, mercury, gold), DMSA (lead, mercury), penicillamine (copper, lead). Complications: nephrotoxicity, electrolyte disturbances, redistribution of metals.

\medterm{Child Poisoning}-- See pediatric toxicology.

\medterm{Cholecystokinin}-- See physiology.

\medterm{Chronic Kidney Disease}-- See nephrology.

\medterm{Chronic Toxicity}-- Adverse effects from repeated or prolonged exposure to low doses. Cumulative damage may be irreversible. Examples: heavy metals (lead, mercury), asbestos, alcohol, solvent exposure. Monitoring: biomarkers, organ function tests.

\medterm{Circle of Willis}-- See neurology.

\medterm{Clearance}-- Volume of plasma cleared of substance per unit time. Renal clearance: urine flow x urine concentration / plasma concentration. Hepatic clearance: blood flow x extraction ratio. Used for dosing adjustments.

\medterm{Clonazepam}-- See benzodiazepines.

\medterm{CO Poisoning}-- See carbon monoxide poisoning.

\medterm{Coma}-- See emergency medicine.

\medterm{Coma Scale}-- See Glasgow Coma Scale.

\medterm{Concentration}-- Amount of substance per unit volume. Units: mg/dL, mmol/L, mEq/L, ppm. Conversion: molecular weight needed. Critical for toxicity assessment.

\medterm{Conjunctivitis}-- See ophthalmology.

\medterm{Contact Dermatitis}-- See dermatology.

\medterm{Controlled Substance}-- See pharmacology.

\medterm{Corrosive}-- Substance causing tissue destruction by chemical action. Acids: coagulation necrosis, eschar. Alkalis: liquefaction necrosis, deep penetration. Ingestion: avoid emesis, endoscopy, surgical consultation.

\medterm{Creatinine}-- See biochemistry.

\medterm{Cross-Tolerance}-- Tolerance to one substance conferring tolerance to another of similar class. Examples: alcohol-benzodiazepines, opioids-methadone.

\medterm{CSF Analysis}-- See neurology.

\medterm{Cumulative Dose}-- Total amount of substance received over time. Important for carcinogens, teratogens, cumulative toxins.

\medterm{Cyanide Poisoning}-- Inhibits cytochrome c oxidase, blocking cellular respiration. Sources: smoke inhalation, industrial exposure, almonds, apricot kernels. Symptoms: headache, confusion, seizures, coma, cardiac arrest, cherry-red skin. Treatment: hydroxocobalamin (first-line), nitrite + thiosulfate kit.

\medterm{Delirium}-- See emergency medicine.

\medterm{Dependence}-- See addiction.

\medterm{Detoxification}-- Removal of toxic substance from body. Methods: activated charcoal, gastric lavage (rarely used), whole bowel irrigation, hemodialysis, hemoperfusion, chelation, supportive care.

\medterm{Diabetic Ketoacidosis}-- See endocrinology.

\medterm{Diffusion}-- Movement of molecules from high to low concentration. Important for drug absorption, toxin distribution.

\medterm{Digitalis (Digoxin) Toxicity}-- Nausea, vomiting, visual changes (yellow halos), arrhythmias. Risk: narrow therapeutic index, renal impairment, hypokalemia. Treatment: digoxin immune Fab, cardiac monitoring, correct electrolytes.

\medterm{Dimercaprol (BAL)}-- Chelating agent for arsenic, mercury, gold. Given IM. Side effects: hypertension, tachycardia, nausea, pain at injection site.

\medterm{Dose}-- Amount of substance administered. Factors: weight, age, renal/hepatic function, drug interactions. Toxic dose varies by substance and individual.

\medterm{Drug Interaction}-- See pharmacology.

\medterm{Drug Screening}-- Urine, blood, hair testing for substances of abuse. Panels: opiates, cocaine, amphetamines, cannabinoids, PCP, benzodiazepines, barbiturates. False positives possible. Confirmation by GC/MS or LC-MS/MS.

\medterm{Drug Testing}-- See drug screening.

\medterm{Drug Therapeutics}-- See pharmacology.

\medterm{Drug Toxicity}-- Adverse effects from drugs at normal or elevated doses. Idiosyncratic reactions, cumulative toxicity, organ damage.

\medterm{Drug Withdrawal}-- See withdrawal syndromes.

\medterm{Dx}-- See diagnosis.

\medterm{Dialysis}-- See nephrology.

\medterm{Edema}-- See general.

\medterm{Effective Dose}-- Dose producing desired effect in 50% of population (ED50). Therapeutic index = TD50/ED50.

\medterm{Elimination Half-Life}-- Time for plasma concentration to decrease by 50%. Determines dosing interval, accumulation, withdrawal from exposure. Affected by renal/hepatic function, age, drug interactions.

\medterm{Emetics}-- See vomiting.

\medterm{Encephalopathy}-- Brain dysfunction from metabolic, toxic, or structural causes. Hepatic, uremic, toxic encephalopathy.

\medterm{Endocrine Disruptor}-- Chemical interfering with hormone systems. Examples: BPA, phthalates, DDT, PCBs. Effects: reproductive disorders, developmental problems, cancer.

\medterm{Endotoxin}-- Lipopolysaccharide from Gram-negative bacteria. Pyrogenic, causes fever, shock, DIC. Contamination concern in pharmaceuticals, medical devices.

\medterm{Enterohepatic Circulation}-- Recycling of toxins between liver and intestine via bile. Prolonged exposure, increased toxicity. Interrupted by activated charcoal, cholestyramine. Examples: morphine, digoxin, estrogen, thyroxine.

\medterm{Epidemiology}-- See public health.

\medterm{Ethylene Glycol Poisoning}-- Antifreeze poisoning. Metabolized to toxic acids (glycolic, oxalic). Causes high anion gap metabolic acidosis, calcium oxalate crystals in urine, renal failure. Treatment: fomepizole (alcohol dehydrogenase inhibitor), hemodialysis, calcium gluconate if hypocalcemia.

\medterm{Ethosuximide}-- See neurology.

\medterm{Ethylene Glycol}-- See ethylene glycol poisoning.

\medterm{Ethosuximide}-- See pharmacology.

\medterm{Excretion}-- Removal of substance from body. Routes: renal (urine), hepatic (bile), pulmonary (gas), sweat, breast milk. Important for toxicity and drug clearance.

\medterm{Exposure Route}-- Path by which substance enters body. Routes: inhalation, ingestion, dermal, injection, ocular. Determines absorption rate and toxicity.

\medterm{Eye Irritation}-- Chemical injury to eye. Immediate irrigation with water or saline for 15-20 minutes. Alkali burns more dangerous than acid. Ophthalmology consultation.

\medterm{Fat-Soluble Toxins}-- Lipophilic toxins stored in fat, prolonged elimination. Examples: organochlorine pesticides, PCBs, dioxins, tetrachlorodiphenylmethane.

\medterm{Federal Register}-- See regulatory affairs.

\medterm{Fentanyl}-- See opioids.

\medterm{Fiber Optic}-- See radiology.

\medterm{Flucloxacillin}-- See pharmacology.

\medterm{Flumazenil}-- Benzodiazepine antagonist. Reverses sedation but may precipitate seizures in chronic users or tricyclic overdose. Dose: 0.2 mg IV, repeat as needed.

\medterm{Food Additive}-- Substance added to food for preservation, flavor, color. Safetyassess by FDA, EFSA. GRAS (Generally Recognized as Safe) status. Concerns: allergies, carcinogenicity, cumulative exposure.

\medterm{Food Poisoning}-- See infectious disease.

\medterm{Foreign Body}-- See medical implants.

\medterm{Formaldehyde}-- Preservative, industrial chemical. Irritant, carcinogen (nasopharyngeal cancer). Sensitizer. Use in fume hood, PPE required.

\medterm{Gastric Emptying}-- See physiology.

\medterm{Gastroenteritis}-- See gastroenterology.

\medterm{Gastrointestinal Decontamination}-- Methods: activated charcoal, gastric lavage (rare), whole bowel irrigation (sorbitol, PEG). Timing critical for charcoal (1-2 hours post-ingestion).

\medterm{Genotoxicity}-- Ability of substance to damage DNA. Tests: Ames test, chromosomal aberration, micronucleus. Predicts carcinogenicity.

\medterm{Glucose-6-Phosphatase}-- See biochemistry.

\medterm{Glutathione}-- Antioxidant tripeptide (Glu-Cys-Gly). Detoxifies reactive metabolites (e.g., NAPQI from acetaminophen). Depletion causes toxicity. N-acetylcysteine replenishes glutathione.

\medterm{Gout}-- See rheumatology.

\medterm{Half-Life}-- See elimination half-life, biological half-life.

\medterm{Hand, Foot, and Mouth Disease}-- See pediatrics.

\medterm{Heavy Metal Poisoning}-- Lead, mercury, arsenic, cadmium, thallium. Chelation therapy, supportive care, prevention. Sources: paint, water pipes, fish, occupational exposure.

\medterm{Hemodialysis}-- See nephrology.

\medterm{Hemoperfusion}-- Blood passed through adsorbent column (charcoal, resin) to remove toxins. Used for drug overdoses with high protein binding, large volume of distribution.

\medterm{Hemorrhagic}-- See hematology.

\medterm{Hepatic Encephalopathy}-- See geriatrics.

\medterm{Hepatotoxicity}-- Liver damage from toxins, drugs, chemicals. Acetaminophen most common cause of acute liver failure. Monitoring: LFTs, INR, clinical assessment.

\medterm{Hyaline Membrane}-- See pediatrics.

\medterm{Hyperemesis}-- See gastroenterology.

\medterm{Hypersensitivity}-- See pharmacology.

\medterm{Hypertensive}-- See cardiology.

\medterm{Hypokalemia}-- Low potassium. Causes: diuretics, vomiting, diarrhea, hyperaldosteronism. Symptoms: weakness, arrhythmias, ECG changes. Treatment: potassium replacement.

\medterm{Idiopathic}-- See general.

\medterm{Idiosyncratic Reaction}-- Unpredictable adverse reaction not related to dose or pharmacology. Immunological or metabolic origin. Examples: drug-induced liver injury, aplastic anemia. Difficult to predict.

\medterm{Immunotoxicity}-- Immune system damage from chemicals. Effects: immunosuppression, hypersensitivity, autoimmunity. Examples: benzene, chemotherapy, biologics.

\medterm{Inhalation Injury}-- See pulmonology.

\medterm{Ingestion}-- See poisoning.

\medterm{Insulin}-- See endocrinology.

\medterm{Interference}-- See pharmacology.

\medterm{Intoxication}-- See poisoning.

\medterm{Ipecac}-- Emetic no longer recommended for routine poisoning management. Risk of aspiration, delays charcoal administration.

\medterm{Isoniazid Toxicity}-- Hepatotoxicity, peripheral neuropathy, seizures. Treatment: pyridoxine (vitamin B6) supplementation, activated charcoal, seizure control.

\medterm{Iatrogenic}-- See general.

\medterm{Kidney Failure}-- See nephrology.

\medterm{Kinetics}-- See pharmacology.

\medterm{Lactate Dehydrogenase}-- See biochemistry.

\medterm{LD50}-- Lethal dose for 50% of test population. Measure of acute toxicity. Lower LD50 = more toxic. Measured in mg/kg body weight. Ethical concerns limit animal testing.

\medterm{Lead Poisoning}-- Neurotoxic, hematological, renal effects. Sources: paint, water, gasoline, occupational. Children more vulnerable. Treatment: chelation (succimer, EDTA, dimercaprol), remove exposure. Blood lead level guides treatment.

\medterm{Liver}-- See biochemistry.

\medterm{Liver Transplant}-- See gastroenterology.

\medterm{Local Anesthetic Toxicity}-- CNS toxicity: tinnitus, perioral numbness, seizures. Cardiovascular toxicity: arrhythmias, hypotension, cardiac arrest. Treatment: lipid emulsion therapy (Intralipid), seizure control, ACLS.

\medterm{Long-Term Effect}-- See chronic toxicity.

\medterm{Loss of Appetite}-- See gastroenterology.

\medterm{Low Molecular Weight Heparin}-- See hematology.

\medterm{Lung Cancer}-- See oncology.

\medterm{Lymphedema}-- See geriatrics.

\medterm{Magnetic Resonance Imaging}-- See radiology.

\medterm{Maltol}-- See food additive.

\medterm{Manifestation}-- See clinical presentation.

\medterm{Manganese}-- Neurotoxic at high exposure. Manganism: Parkinson-like symptoms. Sources: welding, batteries, contaminated water.

\medterm{Mechanism of Toxicity}-- How substance causes harm. Examples: enzyme inhibition, membrane disruption, receptor binding, oxidative stress, DNA damage.

\medterm{Melanoma}-- See dermatology.

\medterm{Mercury Poisoning}-- Three forms: elemental (inhalation, neurological), inorganic (renal, GI), organic (methylmercury, neurological). Minamata disease from methylmercury. Treatment: chelation (dimercaprol, DMSA), remove exposure.

\medterm{Micrograms}-- See units.

\medterm{Micromolar}-- See units.

\medterm{Mild Cognitive Impairment}-- See neurology.

\medterm{Milliequivalent}-- See units.

\medterm{Millimolar}-- See units.

\medterm{Minimum Alveolar Concentration}-- See pharmacology.

\medterm{Minimum Effect Concentration}-- See pharmacology.

\medterm{Minimum Lethal Concentration}-- See pharmacology.

\medterm{Molecular Weight}-- Mass of molecule. Important for dosing, toxicity calculations.

\medterm{Mortality Rate}-- See public health.

\medterm{Movement Disorder}-- See neurology.

\medterm{Mucosal Irritation}-- Chemical injury to mucous membranes. Eyes, nose, throat, GI tract. Treatment: irrigation, supportive care.

\medterm{Mutagen}-- Substance causing genetic mutations. Tests: Ames test, chromosomal aberration. Correlated with carcinogenicity.

\medterm{N-acetylcysteine}-- See acetaminophen toxicity.

\medterm{NAD+}-- See biochemistry.

\medterm{Nephrotoxicity}-- Kidney damage from toxins, drugs, contrast. Aminoglycosides, cisplatin, contrast media, ethylene glycol. Monitoring: creatinine, urine output. Prevention: hydration, dose adjustment.

\medterm{Neurotoxin}-- Substance toxic to nervous system. Examples: lead, mercury, organophosphates, botulinum toxin, tetrodotoxin, snake venoms.

\medterm{Nephropathy}-- See nephrology.

\medterm{Nicotine}-- See tobacco.

\medterm{Night Blindness}-- See ophthalmology.

\medterm{Nomogram}-- Graphical calculation tool. Rumack-Matthew nomogram for acetaminophen toxicity. Guides NAC treatment.

\medterm{Nuclear Medicine}-- See nuclear medicine.

\medterm{Nursing}-- See occupational health.

\medterm{Occupational Exposure}-- Workplace exposure to toxins. Monitoring, PPE, engineering controls, regulations (OSHA, NIOSH).

\medterm{Occupational Medicine}-- See occupational health.

\medterm{Odor Threshold}-- Lowest concentration detectable by smell. Important for warning properties of toxins.

\medterm{Odynophagia}-- See gastroenterology.

\medterm{Oliguria}-- See nephrology.

\medterm{Oncogenic}-- Cancer-causing. See carcinogen.

\medterm{Operative}-- See surgery.

\medterm{Opportunistic Infection}-- See infectious disease.

\medterm{Organophosphate Poisoning}-- Acetylcholinesterase inhibitors. SLUDGE syndrome: salivation, lacrimation, urination, defecation, GI upset, emesis. Miosis, bradycardia, bronchospasm, muscle fasciculations. Treatment: atropine, pralidoxime, supportive care.

\medterm{Opioid Antagonist}-- See naloxone.

\medterm{Oral Toxicity}-- See ingestion.

\medterm{Organophosphate}-- See organophosphate poisoning.

\medterm{Ovarian Cancer}-- See oncology.

\medterm{Overdose}-- See poisoning.

\medterm{Oxidative Stress}-- See biochemistry.

\medterm{Paraquat}-- Herbicide, highly toxic if ingested. Pulmonary fibrosis, renal failure. No antidote. Poor prognosis.

\medterm{Paracetamol}-- See acetaminophen.

\medterm{Particulate Matter}-- See air pollution.

\medterm{Pathogenesis}-- See pathophysiology.

\medterm{Pathology}-- See pathology.

\medterm{Pediatric Toxicology}-- Poisoning in children. Common: iron, medications, household cleaners, plants. Prevention: child-proof containers, education.

\medterm{Penicillin}-- See biochemistry.

\medterm{Periodontal Disease}-- See dentistry.

\medterm{Persistent Organic Pollutant}-- POCPs: organochlorines, PCBs, dioxins. Bioaccumulate, long-range transport, persistent. Stockholm Convention bans/restricts.

\medterm{Pharmaceutical}-- See pharmacology.

\medterm{Pharmacokinetics}-- See pharmacology.

\medterm{Pharmacology}-- See pharmacology.

\medterm{Pharmacodynamics}-- See pharmacology.

\medterm{Pharmacovigilance}-- Monitoring drug safety after approval. Adverse event reporting, signal detection, risk management.

\medterm{Phencyclidine (PCP)}-- See hallucinogens.

\medterm{Photoallergy}-- See dermatology.

\medterm{Phototoxicity}-- Skin reaction from UV light + chemical. Sunburn-like. Examples: tetracyclines, thiazides, NSAIDs, psoralens.

\medterm{Plasma Concentration}-- See pharmacokinetics.

\medterm{Poison Control}-- See emergency medicine.

\medterm{Poisoning}-- Harmful effect from exposure to toxic substance. Routes: ingestion, inhalation, dermal, injection. Management: ABCs, decontamination, antidote if available, supportive care.

\medterm{Population}-- See public health.

\medterm{Potency}-- See pharmacology.

\medterm{Prenatal Exposure}-- See teratogen.

\medterm{Prescription Drug}-- See pharmacology.

\medterm{Prevention}-- See toxicology.

\medterm{Primary Metabolism}-- See hepatic metabolism.

\medterm{Probable Cause}-- See diagnosis.

\medterm{Procedure}-- See surgery.

\medterm{Product Liability}-- See regulatory affairs.

\medterm{Prognosis}-- See diagnosis.

\medterm{Prophylaxis}-- See prevention.

\medterm{Protein Binding}-- Drug bound to plasma proteins (mainly albumin). Only free drug is active. Displacement interactions possible. Low protein binding = easier dialyzable.

\medterm{Proteinuria}-- See nephrology.

\medterm{Proton Pump Inhibitor}-- See gastroenterology.

\medterm{Proximate Cause}-- See legal toxicology.

\medterm{Psychoactive Substance}-- See drugs of abuse.

\medterm{Psychosis}-- See psychiatry.

\medterm{Pulmonary Edema}-- See pulmonology.

\medterm{Pulse Oximetry}-- See emergency medicine.

\medterm{Purging}-- See eating disorders.

\medterm{Pyridoxine}-- Vitamin B6. Antidote for isoniazid toxicity, prevents peripheral neuropathy.

\medterm{Quality Control}-- See regulatory affairs.

\medterm{Radiation Poisoning}-- See radiation injury.

\medterm{Refeeding Syndrome}-- See clinical nutrition.

\medterm{Referred Pain}-- See pain.

\medterm{Relative Risk}-- See epidemiology.

\medterm{Renal Replacement Therapy}-- See nephrology.

\medterm{Reproductive Toxicity}-- Effects on reproduction: infertility, teratogenesis, fetal toxicity, developmental toxicity. Testing required for pharmaceuticals.

\medterm{Respiratory Depression}-- See emergency medicine.

\medterm{Respiratory Failure}-- See pulmonology.

\medterm{Respiratory Protection}-- See occupational health.

\medterm{Retention Time}-- See pharmacokinetics.

\medterm{Retinopathy}-- See ophthalmology.

\medterm{Reverse Transcription}-- See molecular biology.

\medterm{Risk Assessment}-- Process: hazard identification, dose-response, exposure assessment, risk characterization. Used for regulatory decisions.

\medterm{Risk-Benefit Analysis}-- Weighing therapeutic benefit against potential harm.

\medterm{Risk Factor}-- See epidemiology.

\medterm{Rocuronium}-- See pharmacology.

\medterm{Running Total}-- See exposure.

\medterm{Salicylate}-- See aspirin toxicity.

\medterm{Saturation}-- See pharmacokinetics.

\medterm{Schedule I}-- See controlled substances.

\medterm{Schedule II}-- See controlled substances.

\medterm{Schedule III}-- See controlled substances.

\medterm{Schedule IV}-- See controlled substances.

\medterm{Schedule V}-- See controlled substances.

\medterm{Seizure}-- See neurology.

\medterm{Sedative-Hypnotic}-- CNS depressants. Benzodiazepines, barbiturates, Z-drugs. Overdose: respiratory depression, coma. Treatment: supportive, flumazenil (benzodiazepines only), airway management.

\medterm{Second-Hand Smoke}-- See tobacco.

\medterm{Secretion}-- See pharmacokinetics.

\medterm{Self-Experimentation}-- See research ethics.

\medterm{Serum Concentration}-- See pharmacokinetics.

\medterm{Serum Sickness}-- Type III hypersensitivity reaction. Fever, rash, arthritis, proteinuria. Caused by drugs (penicillin, cephalosporins), antivenoms. Treatment: stop offending agent, antihistamines, corticosteroids.

\medterm{Sexual Dysfunction}-- See pharmacology.

\medterm{Side Effect}-- See adverse effect.

\medterm{Simethicone}-- See gastroenterology.

\medterm{Single-Dose Toxicity}-- See acute toxicity.

\medterm{Skin Absorption}-- See dermal exposure.

\medterm{Skin Irritation}-- See dermatology.

\medterm{Social Isolation}-- See psychiatry.

\medterm{Social Support}-- See psychiatry.

\medterm{Soft Tissue Injury}-- See emergency medicine.

\medterm{Solvent Exposure}-- See occupational medicine.

\medterm{Source Control}-- See infection.

\medterm{Sodium Polystyrene Sulfonate}-- Potassium-binding resin for hyperkalemia. Risk: colonic necrosis, especially with sorbitol.

\medterm{Soil Contamination}-- Environmental toxicology. Heavy metals, pesticides, industrial waste. Remediation: excavation, bioremediation, containment.

\medterm{Specificity}-- See pharmacology.

\medterm{Spinal Cord Injury}-- See neurology.

\medterm{Spill}-- See emergency response.

\medterm{St. John's Wort}-- Herbal supplement. CYP3A4 inducer, reduces efficacy of many drugs (oral contraceptives, warfarin, antiretrovirals).

\medterm{Steady State}-- See pharmacokinetics.

\medterm{Substance Abuse}-- See addiction.

\medterm{Substance Dependence}-- See addiction.

\medterm{Substance Use Disorder}-- See addiction.

\medterm{Sulfa Drugs}-- See sulfonamides.

\medterm{Sulfonamide}-- See pharmacology.

\medterm{Supportive Care}-- Mainstay of poison management. Airway, breathing, circulation, IV access, monitoring, treat complications.

\medterm{Suspected Cause}-- See diagnosis.

\medterm{Syncope}-- See cardiology.

\medterm{Systemic Effects}-- Whole-body effects of toxins. Cardiac, neurological, renal, hepatic, hematological.

\medterm{Teratogen}-- Substance causing birth defects. Examples: thalidomide, isotretinoin, alcohol, warfarin, valproate, ACE inhibitors. Prenatal exposure critical.

\medterm{Testicular Cancer}-- See oncology.

\medterm{Therapeutic Index}-- Ratio of toxic dose to therapeutic dose. TI = TD50/ED50. Wide TI = safer drug. Narrow TI drugs require monitoring (digoxin, lithium, warfarin, phenytoin).

\medterm{Therapeutic Window}-- See therapeutic index.

\medterm{Thiamine}-- See biochemistry.

\medterm{Therapy}-- See treatment.

\medterm{Thrombocytopenia}-- See hematology.

\medterm{Tidal Volume}-- See physiology.

\medterm{Tissue Distribution}-- See pharmacokinetics.

\medterm{Tobacco}-- See nicotine.

\medterm{Tolerance}-- See pharmacology.

\medterm{Toxic Dose}-- Dose causing toxicity. Varies by substance, individual, route, duration.

\medterm{Toxicology}-- Study of poisons and their effects. Branches: clinical, forensic, environmental, occupational, regulatory.

\medterm{Toxin}-- Poison produced by living organism. Bacterial (botulinum, tetanus, diphtheria), animal (snake venom, spider venom, scorpion venom), plant (ricin, amanitin), fungal (aflatoxin).

\medterm{Trace Elements}-- Minerals needed in small amounts. Zinc, copper, manganese, selenium, chromium, molybdenum, iodine. Deficiency or toxicity possible.

\medterm{Treatment}-- See management.

\medterm{Tricyclic Antidepressant Overdose}-- Anticholinergic effects, cardiac toxicity (QRS prolongation, arrhythmias), seizures, hypotension. Treatment: sodium bicarbonate (QRS >100 ms), seizure control, supportive care.

\medterm{Tylenol}-- See acetaminophen.

\medterm{Typhoid}-- See infectious disease.

\medterm{Tyramine}-- See pharmacology.

\medterm{Uremia}-- See nephrology.

\medterm{Urine Analysis}-- See nephrology.

\medterm{Urine Toxicology}-- See drug screening.

\medterm{UV Light}-- See phototoxicity.

\medterm{Vaccine}-- See pharmacology.

\medterm{Valproate}-- See pharmacology.

\medterm{Vascular Access}-- See emergency medicine.

\medterm{Venom}-- See toxin.

\medterm{Venomous Bite}-- See emergency medicine.

\medterm{Vomiting}-- Reflex expulsion of gastric contents. Causes: GI, neurological, metabolic, toxic, psychological. Treatment: antiemetics, identify cause, prevent aspiration.

\medterm{Volume of Distribution}-- Vd = dose / plasma concentration. Indicates tissue distribution. High Vd = extensive tissue binding, harder to remove by dialysis.

\medterm{Warfarin}-- See pharmacology.

\medterm{Water Soluble Toxins}-- Excreted by kidneys. Dialysis effective for small molecules with low protein binding (ethanol, methanol, ethylene glycol, lithium, salicylates).

\medterm{Withdrawal Syndrome}-- See withdrawal syndromes.

\medterm{Xanax}-- See benzodiazepines.

\medterm{Xenobiotic}-- Foreign chemical substance in body. Drugs, pollutants, pesticides, industrial chemicals. Metabolized by liver enzymes.

\medterm{Zolpidem}-- See sedative-hypnotics.

\medterm{Zinc}-- Trace element. Deficiency: growth retardation, delayed wound healing, hypogonadism, dysgeusia. Sources: meat, shellfish, legumes, nuts.
